\documentclass[aps,prl,reprint,superscriptaddress,nofootinbib]{revtex4-2}

\usepackage{amsmath,amssymb,amsfonts,bm,graphicx}
\usepackage{hyperref}
\hypersetup{colorlinks=true,linkcolor=blue,citecolor=blue,urlcolor=blue}

\newcommand{\ket}[1]{|#1\rangle}
\newcommand{\bra}[1]{\langle#1|}

\begin{document}
Before evaluating this on the hyperboloid it pays to run the same family through the flat case, where the answer is known and sharp. For a single mode the Fock blocks are elementary~, and $\cos\theta\ket{0}+\sin\theta\ket{1}$ has the Wigner function $(2/\pi)e^{-2|\alpha|^2}f(|\alpha|)$ with
\begin{equation}
f(t)=\cos2\theta+4\sin^2\theta\,t^2-4\sin\theta\cos\theta\,t.
\label{eq:wigner}
\end{equation}
The discriminant of Eq.~ is $16\sin^4\theta$, which is positive for every $\theta\neq0$, so $f$ has real roots and the Wigner function goes negative. The minimum value is $-\sin^2\theta$ and it sits at $|\alpha|=\cot\theta/2$. Hudson's theorem appears here in a concrete form. The admissible window is exactly a point, the negativity retreats to infinity as $\theta\to0$ rather than disappearing, and it shrinks only quadratically in the mixing angle.

The runaway is the mechanism worth keeping in mind. In the flat case the two ratios $W_{11}/W_{00}=4|\alpha|^2-1$ and $|C|/W_{00}=2|\alpha|$ both diverge, so the band of forbidden angles slides towards zero and however small $\theta$ is there is always a radius far enough out to catch it. Section~ shows the hyperboloid closes that escape route.

Minimising Eq.~ over $\phi$ is immediate. Giving the excited amplitude a relative phase $e^{i\delta}$ replaces $\cos\phi$ by $\cos(\phi+\delta)$ and so changes nothing after that minimisation. Hence $W_\theta\ge0$ on the whole hyperboloid if and only if
\begin{multline}
\cos^2\theta\,W_{00}(\tau)+\sin^2\theta\,W_{11}(\tau)\\
-2\,[\sin\theta\cos\theta]\,|C(\tau)|\;\ge\;0
\label{eq:positive}
\end{multline}
for every $\tau\ge0$. Positivity is therefore a one-parameter question at each $\tau$, and the admissible $\theta$ is bounded by the smaller root of a quadratic.

The three functions play distinct roles, and Fig.~(a) shows them. $W_{00}$ is positive everywhere and supplies the budget. $W_{11}$ changes sign and is negative near the origin. The coherence $C$ enters through $\cos\phi$, so whatever its sign it subtracts somewhere on each orbit of the rotation generated by $\hat K_0$.

\begin{figure}[t]
\centering
\includegraphics[width=0.95\columnwidth]{fig1.pdf}
\caption{(a) The three radial blocks at $k=1$. $W_{00}$ stays positive, $W_{11}$ is negative out to $\tau\simeq0.67$, and the coherence $|C|$ peaks at $\tau\simeq0.60$. (b) The positivity criterion~, normalised by $W_{00}$, for five mixing angles. The curve at $\theta_W$ touches zero at $\tau=1.225$, away from both the origin and the tail, and the curve at $30^\circ$ dips below. Normalising by $W_{00}$ removes the common exponential decay and makes the margin quoted in the text visible.}
\label{fig:blocks}
\end{figure}

Evaluating Eq.~ at $k=1$ gives
\begin{equation}
\theta_W=24.92^\circ,\qquad \sin^2\theta_W=0.178,
\label{eq:thetaW}
\end{equation}
as the largest admissible mixing angle. The Wigner-positive pure states contain the Perelemov coherent orbit strictly, and Hudson's dichotomy has no analogue here.

Comparing with Eq.~ makes the content of Eq.~ plain. The flat family is negative for every $\theta>0$ and the hyperbolic family is positive up to a finite angle, so the two definitions of classicality that Hudson's theorem identities come apart as soon as the phase space is curved. The gap is not a boundary effect either. At $\theta=20^\circ$ the state carries $12\%$ of its weight on the first excited level and its Wigner function stays positive across the whole hyperboloid.

The binding value of $\tau$ is $1.225$ at $k=1$ [Fig.~(b)]. It sits away from the origin, where $W_{11}$ is deepest but $W_{00}$ is also largest, and away from the tail, where Sec.~ shows the quadratic form becomes marginal. The threshold is therefore fixed at intermediate hyperbolic distance, and the margin elsewhere stays finite.
\end{document}
